\documentclass[11pt]{article}
\usepackage[a4paper,margin=1in]{geometry}
\usepackage{fontspec}
\usepackage[UTF8]{ctex}
\setmainfont{TeX Gyre Termes}
\usepackage{amsmath,amssymb,mathtools,bm}
\usepackage{xcolor}
\usepackage{booktabs}
\usepackage{enumitem}
\usepackage{hyperref}
\usepackage{tcolorbox}
\hypersetup{colorlinks=true,linkcolor=blue,urlcolor=blue,citecolor=blue}
\setlength{\parindent}{0pt}
\setlength{\parskip}{0.35em}
\newtcolorbox{formulabox}[1]{
  colback=blue!3!white,
  colframe=blue!55!black,
  title=#1,
  fonttitle=\bfseries,
  boxsep=1mm,
  left=1mm,right=1mm,top=1mm,bottom=1mm
}
\newtcolorbox{insightbox}[1]{
  colback=orange!4!white,
  colframe=orange!55!black,
  title=#1,
  fonttitle=\bfseries,
  boxsep=1mm,
  left=1mm,right=1mm,top=1mm,bottom=1mm
}
\newtcolorbox{derivationbox}[1]{
  colback=green!3!white,
  colframe=green!45!black,
  title=#1,
  fonttitle=\bfseries,
  boxsep=1mm,
  left=1mm,right=1mm,top=1mm,bottom=1mm
}
\newcommand{\E}{\mathbb{E}}
\newcommand{\Var}{\mathrm{Var}}
\newcommand{\Cov}{\mathrm{Cov}}
\newcommand{\KL}{\mathrm{KL}}
\newcommand{\MI}{\mathrm{I}}

\title{Chapter 5 Notes\\\large Classification}
\author{AI for Science}
\date{\today}

\begin{document}
\maketitle
\tableofcontents
\newpage

\section{导读：分类不只是画一条边界}
分类（Classification）研究的是如何把输入 $\mathbf x$ 映射到离散类别集合。要真正讲清这一章，必须把“判别函数”“后验概率”“决策规则”“损失函数”这些概念接成一条完整链条。主线是
\[
\text{判别函数} \rightarrow \text{Bayes classifier} \rightarrow \text{风险最小化} \rightarrow \text{生成式与判别式} \rightarrow \text{Logistic / Softmax}.
\]

\begin{insightbox}{学习目标}
\begin{itemize}[leftmargin=1.6em]
  \item 掌握线性判别边界的几何含义。
  \item 理解最小错误率决策和代价敏感决策的区别。
  \item 掌握 logistic / softmax 模型的概率解释与交叉熵推导。
\end{itemize}
\end{insightbox}

\section{5.1 判别函数与线性决策边界}
\begin{formulabox}{二分类线性判别函数}
\begin{equation}
y(\mathbf x)=\mathbf w^\top \mathbf x+w_0,
\qquad
\hat c=\begin{cases}C_1,&y(\mathbf x)\ge 0,\\ C_2,&y(\mathbf x)<0.
\end{cases}
\end{equation}
\textbf{变量释义：} $\mathbf w$ 是法向量，$w_0$ 决定超平面偏移。\\
\textbf{推导思路（2-4步）：}
\begin{enumerate}[leftmargin=1.6em]
  \item 判别函数输出一个实数打分。
  \item 打分为零的点组成决策边界。
  \item 边界两侧对应不同类别。
\end{enumerate}
\textbf{何时使用：} 类别大致线性可分时的基础分类器。
\end{formulabox}

\begin{formulabox}{多分类线性判别}
\begin{equation}
\hat c=\arg\max_k y_k(\mathbf x),
\qquad y_k(\mathbf x)=\mathbf w_k^\top\mathbf x+w_{k0}.
\end{equation}
\textbf{变量释义：} 每个类别都有一个打分函数，取最大得分对应的类别。\\
\textbf{推导思路（2-4步）：}
\begin{enumerate}[leftmargin=1.6em]
  \item 多分类时不再只比较正负，而是比较各类得分大小。
  \item 两个类别边界由 $y_i(\mathbf x)=y_j(\mathbf x)$ 给出。
  \item 因而多分类线性模型仍对应若干超平面拼接出来的区域划分。
\end{enumerate}
\textbf{何时使用：} softmax 之前的 logits 层、线性多分类基线。
\end{formulabox}

\section{5.2 Bayes classifier 与风险最小化}
\begin{formulabox}{最小错误率分类}
\begin{equation}
\hat c(\mathbf x)=\arg\max_k p(C_k\mid \mathbf x).
\end{equation}
\textbf{变量释义：} $p(C_k\mid \mathbf x)$ 是给定输入后的类别后验概率。\\
\textbf{推导思路（2-4步）：}
\begin{enumerate}[leftmargin=1.6em]
  \item 在 0-1 损失下，预测错记 1，预测对记 0。
  \item 每个动作的条件风险就是预测错误的概率。
  \item 因而选择后验概率最大的类别即可使错误率最小。
\end{enumerate}
\textbf{何时使用：} 默认分类决策规则与理论最优分类器分析。
\end{formulabox}

\begin{formulabox}{代价敏感决策}
\begin{equation}
R(\alpha_i\mid \mathbf x)=\sum_{k=1}^{K} \lambda_{ik}p(C_k\mid \mathbf x),
\qquad
\hat\alpha=\arg\min_i R(\alpha_i\mid \mathbf x).
\end{equation}
\textbf{变量释义：} $\lambda_{ik}$ 表示真实类别为 $C_k$ 时采取动作 $\alpha_i$ 的代价。\\
\textbf{推导思路（2-4步）：}
\begin{enumerate}[leftmargin=1.6em]
  \item 实际任务中不同错误的后果通常不一样。
  \item 因此不能只看谁的后验概率大，还要看错误代价大小。
  \item 最优动作是最小化条件期望风险，而非单纯最大后验类别。
\end{enumerate}
\textbf{何时使用：} 医学诊断、欺诈检测、工业告警等高代价不对称场景。
\end{formulabox}

\section{5.3 生成式与判别式}
\begin{formulabox}{两种建模路线}
\begin{align}
\text{生成式：}
&&p(C_k\mid \mathbf x) &\propto p(\mathbf x\mid C_k)p(C_k),\\
\text{判别式：}
&&\text{直接建模 } p(C_k\mid \mathbf x) \text{ 或判别函数 } y_k(\mathbf x).
\end{align}
\textbf{变量释义：} 生成式先描述“各类数据长什么样”，判别式则直接描述“给定输入属于哪一类”。\\
\textbf{推导思路（2-4步）：}
\begin{enumerate}[leftmargin=1.6em]
  \item 生成式模型先建联合分布或类条件分布。
  \item 判别式模型直接针对分类边界或后验概率建模。
  \item 前者更完整，后者通常更贴近预测目标。
\end{enumerate}
\textbf{何时使用：} 小样本、缺失值处理偏向生成式；纯预测精度常偏向判别式。
\end{formulabox}

\section{5.4 Logistic 回归}
\begin{formulabox}{Sigmoid 概率输出}
\begin{align}
p(C_1\mid \mathbf x)&=\sigma(a),
\qquad a=\mathbf w^\top\mathbf x+w_0,\\
\sigma(a)&=\frac{1}{1+e^{-a}}.
\end{align}
\textbf{变量释义：} $a$ 是 logit，$\sigma(a)$ 把任意实数压缩到 $(0,1)$。\\
\textbf{推导思路（2-4步）：}
\begin{enumerate}[leftmargin=1.6em]
  \item 先用线性函数得到实数打分。
  \item 再用 sigmoid 把打分映射为概率。
  \item 于是模型输出可解释为二分类后验概率。
\end{enumerate}
\textbf{何时使用：} 二分类任务、概率预测与可解释线性分类。
\end{formulabox}

\begin{formulabox}{后验 odds 与 logit 的线性关系}
\begin{equation}
\log\frac{p(C_1\mid \mathbf x)}{p(C_2\mid \mathbf x)} = \mathbf w^\top\mathbf x + w_0.
\end{equation}
\textbf{变量释义：} 左边是 posterior odds 的对数，也叫 logit。\\
\textbf{推导思路（2-4步）：}
\begin{enumerate}[leftmargin=1.6em]
\item Logistic 回归直接建模正类后验概率。
\item 把 sigmoid 公式两边取 odds 再取对数，就得到线性形式。
\item 这说明 Logistic 回归实际上是在建模“后验 odds 的对数是输入的线性函数”。
\end{enumerate}
\textbf{何时使用：} 解释 Logistic 回归为何既是概率模型又是线性分类器。
\end{formulabox}
\begin{formulabox}{Bernoulli 似然与交叉熵}
\begin{align}
p(\mathbf t\mid X,\mathbf w)&=\prod_{n=1}^{N} y_n^{t_n}(1-y_n)^{1-t_n},\\
E(\mathbf w)&=-\sum_{n=1}^{N}\left[t_n\log y_n+(1-t_n)\log(1-y_n)\right].
\end{align}
\textbf{变量释义：} $t_n\in\{0,1\}$，$y_n=p(C_1\mid\mathbf x_n)$。\\
\textbf{推导思路（2-4步）：}
\begin{enumerate}[leftmargin=1.6em]
  \item 把每个标签看成 Bernoulli 观测。
  \item 独立样本下，总似然是所有 Bernoulli 概率的乘积。
  \item 取负对数后得到二分类交叉熵。
\end{enumerate}
\textbf{何时使用：} 二分类神经网络输出层和广义线性模型。
\end{formulabox}

\begin{derivationbox}{Logistic 梯度为什么会化简成 $y-t$}
对单个样本，损失为
\[
E_n=-t_n\log y_n-(1-t_n)\log(1-y_n),
\qquad y_n=\sigma(a_n).
\]
先对 $y_n$ 求导，再乘上 $dy_n/da_n=y_n(1-y_n)$，可得
\[
\frac{\partial E_n}{\partial a_n}=y_n-t_n.
\]
这个结果非常重要，因为输出层误差正好就是“预测概率减真实标签”。
\end{derivationbox}

\section{5.5 Softmax 回归}
\begin{formulabox}{Softmax 概率}
\begin{equation}
p(C_k\mid \mathbf x)=\frac{\exp(a_k)}{\sum_{j=1}^{K}\exp(a_j)},
\qquad a_k=\mathbf w_k^\top\mathbf x+w_{k0}.
\end{equation}
\textbf{变量释义：} $a_k$ 是第 $k$ 类的 logit。\\
\textbf{推导思路（2-4步）：}
\begin{enumerate}[leftmargin=1.6em]
  \item 每个类别先产生一个实数打分。
  \item 指数化后变成正权重。
  \item 再除以总和，得到和为 1 的概率分布。
\end{enumerate}
\textbf{何时使用：} 单标签多分类任务的标准输出层。
\end{formulabox}

\begin{formulabox}{Categorical 似然与多分类交叉熵}
\begin{align}
p(T\mid X,\mathbf W)&=\prod_{n=1}^{N}\prod_{k=1}^{K} y_{nk}^{t_{nk}},\\
E(\mathbf W)&=-\sum_{n=1}^{N}\sum_{k=1}^{K} t_{nk}\log y_{nk}.
\end{align}
\textbf{变量释义：} $t_{nk}$ 是 one-hot 标签。\\
\textbf{推导思路（2-4步）：}
\begin{enumerate}[leftmargin=1.6em]
  \item 单个多分类标签服从 categorical 分布。
  \item one-hot 标签让真实类别对应的概率项被保留。
  \item 取负对数后得到多分类交叉熵目标。
\end{enumerate}
\textbf{何时使用：} 多分类网络训练与概率输出解释。
\end{formulabox}

\begin{formulabox}{Softmax 输出层梯度}
\begin{equation}
\frac{\partial E_n}{\partial a_{nk}} = y_{nk}-t_{nk}.
\end{equation}
\textbf{变量释义：} $a_{nk}$ 是第 $n$ 个样本对第 $k$ 类的 logit，$y_{nk}$ 是 softmax 概率。\\
\textbf{推导思路（2-4步）：}
\begin{enumerate}[leftmargin=1.6em]
\item 先对多分类交叉熵对 $y_{nk}$ 求导。
\item 再结合 softmax 的 Jacobian 做链式法则展开。
\item 最终同样化简成“预测概率减 one-hot 标签”。
\end{enumerate}
\textbf{何时使用：} 多分类网络输出层反向传播。
\end{formulabox}
\section{总结与常见误区}
\begin{itemize}[leftmargin=1.6em]
  \item 分类应理解为“概率估计 + 决策规则”，而不是只画一条边界。
  \item Bayes classifier 给出最小错误率决策，代价敏感任务则要最小化期望风险。
  \item Logistic 和 Softmax 的交叉熵都来自明确的概率似然。
  \item 误区1：把最大后验和最小风险混为一谈。
  \item 误区2：认为生成式和判别式只能二选一；它们服务的目标不同。
  \item 误区3：把 softmax 概率直接当作绝对可信置信度，而不做校准检查。
\end{itemize}

\end{document}



